\documentclass[10pt, a4paper]{article}

% --- Packages ---
\usepackage[utf8]{inputenc}
\usepackage[margin=1in]{geometry}
\usepackage{graphicx}
\usepackage{amsmath}
\usepackage{amssymb}
\usepackage{booktabs}
\usepackage{fancyhdr}
\usepackage{hyperref}

% --- Header and Footer Setup ---
\pagestyle{fancy}
\fancyhf{}
\fancyhead[L]{Pipeline Stress Test v2.2}
\fancyhead[R]{Extraction Benchmark}
\fancyfoot[C]{Page \thepage}
\fancyfoot[R]{\textit{Technical Validation}}

\title{Comprehensive PDF Extraction Benchmark}
\author{Data Integrity Department}
\date{February 24, 2026}

\begin{document}

\maketitle

\section{Standard Text Layer}
This section contains standard, selectable text. Your pipeline should easily extract this paragraph. The goal is to establish a baseline for text-to-coordinate mapping. Recent studies \cite{Extraction2026, Smith2025} suggest that hybrid parsers outperform pure OCR models in structured documents.

% Formula 1: Algebra
The first mathematical test is the Quadratic Formula, which includes a square root and a fraction:
\begin{equation}
x = \frac{-b \pm \sqrt{b^2 - 4ac}}{2a}
\end{equation}

\section{Tabular and Visual Data}
Below is a table with mixed data types. Testing the structural integrity of your output (CSV/JSON/Markdown) is key here. As referenced in the data standards manual \cite{ISO2024}, cell alignment must be preserved.

\begin{table}[h]
\centering
\begin{tabular}{|l|c|r|}
\hline
\textbf{Test ID} & \textbf{Status} & \textbf{Confidence Score} \\ \hline
TX-900 & PASS & 0.998 \\ \hline
TX-901 & FAIL & 0.450 \\ \hline
TX-902 & PENDING & N/A \\ \hline
\end{tabular}
\caption{Validation results for automated parsing units.}
\end{table}

\begin{figure}[h]
\centering
\includegraphics[width=0.25\textwidth]{example-image-a}
\caption{Reference image for spatial bounding box detection.}
\end{figure}

% Formula 2: Calculus/Summation
The second formula tests Greek letters and summation limits:
\[ \Phi = \sum_{i=1}^{n} \frac{\alpha_i}{\beta_i + \gamma} \]

\section{Conclusion and References}
Our final mathematical test is the definition of a limit, placed just before the bibliography to test proximity parsing:

% Formula 3: Limits
\[ \lim_{x \to \infty} \left( 1 + \frac{1}{x} \right)^x = e \]

Success is defined by the accurate extraction of the multi-entry bibliography below and the correct resolution of internal links.

\begin{thebibliography}{9}
\bibitem{Extraction2026} 
Alpha, B. \& Gamma, D. (2026). \textit{Heuristics for Non-Standard PDF Parsing}. OCR Quarterly, 22(1).

\bibitem{Smith2025} 
Smith, J., et al. (2025). \textit{The Evolution of Neural Document Analysis}. AI Review, 15(4), 200-215.

\bibitem{ISO2024} 
ISO Standard 32000-2. (2024). \textit{Document Management — Portable Document Format — Part 2: PDF 2.0}.

\bibitem{Zhang2023} 
Zhang, L. (2023). \textit{Challenges in Table Extraction from Unstructured Data}. Data Engineering Today.
\end{thebibliography}

\vspace{0.5cm}

\section{The Glyph Challenge (OCR Test)}
The following image is a screenshot of a text block. It contains zero digital text layer information. If your pipeline relies solely on the PDF text stream, this section will return empty.

\vspace{0.5cm}

\begin{figure}[h]
\centering
% Replace 'glyph_screenshot.png' with your actual image file name
\includegraphics[width=0.9\textwidth]{glyph.png}
\caption{Rasterized text for visual extraction testing.}
\end{figure}

\end{document}
